\documentclass[10pt,a4paper]{amsart}
\usepackage[margin=19mm]{geometry}
\usepackage{amsmath,amssymb,mathtools}
\usepackage[scr=rsfs,cal=euler]{mathalfa}
\usepackage{xfrac}
\usepackage[unicode,hidelinks,bookmarks=false]{hyperref}
\newcommand{\ie}{i.e.\ }
\newcommand{\cf}{cf.\ }
\newcommand{\resp}{respectively\ }
\newcommand{\Subsection}{\subsection}
\providecommand{\nobreakdash}{\nobreak\hbox{-}}
\setlength{\emergencystretch}{2em}
\multlinegap0pt
\usepackage{bm}
\usepackage{bbm}
\usepackage{dsfont}
\usepackage{xspace}
\usepackage{cancel}
\usepackage{mathtools}
\usepackage{cleveref}
\usepackage{amsthm}
\makeatletter
\@ifundefined{theorem}{\newtheorem{theorem}{Theorem}}{}
\makeatother
\DeclareMathOperator{\Aut}{Aut}
\DeclareMathOperator{\aind}{aind}
\DeclareMathOperator{\eind}{eind}
\DeclareMathOperator{\gr}{gr}
\DeclareMathOperator{\ind}{ind}
\DeclareMathOperator{\ldg}{ldg}
\DeclareMathOperator{\sgn}{sgn}
\newcommand{\NN}{\mathbb{N}}
\newcommand{\abs}[1]{\left\lvert#1\right\rvert}
\title[Edge inducibility via local directed graphs]{Edge inducibility via local directed graphs}
\author{Ting-Wei Chao, Asaf Cohen Antonir, Anqi Li, Hung-Hsun Hans Yu}
\date{Source: arXiv:2509.24064v1 (CC BY 4.0)}
\begin{document}
\maketitle

\noindent\emph{Source and licence.} Edge inducibility via local directed graphs. Version arXiv:2509.24064v1 (CC BY 4.0), distributed under the Creative Commons Attribution 4.0 International licence, as recorded in the arXiv licence metadata for this exact version and shown on the abstract page https://arxiv.org/abs/2509.24064v1. Full rights notice: licenses/arxiv-2509.24064-CC-BY-4.0.txt.\par\smallskip

\noindent\textit{Editorial scope.} Counted unit: the Theorem whose source label is \texttt{thm:aind-eind-ind} in the article (source0.tex lines 1328--1332, source label \texttt{thm:aind-eind-ind}), delivered together with the article's own complete original proof of it (source0.tex lines 1333--1372, one \texttt{proof} environment of 40 lines). No mathematics is written, repaired or re-derived: statement and proof are the article's own text line for line. Required context retained uncounted: thm:eind-as-hard-as-ind (lines 718--721). Every definition, display and label of those lines is retained unchanged. Omitted from this package and not redistributed: 1--717, 722--1327, 1373--1973. Nothing inside a retained excerpt is omitted or altered: every retained body is delivered exactly as the article's lines are written, so no omission receipt of any kind accompanies this package. Citations: the retained excerpts carry no citation command at all, so no bibliography entry of the article is transcribed and no reference is redistributed. Reference adaptation: none was needed and none was made; every reference inside the delivered unit resolves inside it, so no unresolved reference or citation is produced. Printed slips kept literal: no printed word, symbol, hypothesis or number of the counted statement or of its proof was altered. Layout and class: the article's own class is \texttt{amsart} with packages including amsthm, amsmath, amssymb, stmaryrd, inputenc, fontenc, tikz, ifthen, array, booktabs, systeme, enumitem, microtype, bm; the wrapper is the lane's pinned \texttt{amsart} editorial wrapper at 10pt on A4, which restates the theorem-like environments the retained text needs and adds the article's own macro definitions that the retained text needs (\texttt{\textbackslash Aut}, \texttt{\textbackslash NN}, \texttt{\textbackslash abs}, \texttt{\textbackslash aind}, \texttt{\textbackslash eind}, \texttt{\textbackslash gr}, \texttt{\textbackslash ind}, \texttt{\textbackslash ldg}, \texttt{\textbackslash sgn}), with extra wrapper packages (bm, bbm, dsfont, xspace, cancel, mathtools, cleveref). The delivered numbering is the unit's own. Assets: no retained span contains a figure environment, a graphics-inclusion command, a PDF-page inclusion, a table or a TikZ picture; no image file is copied into this package, the package declares no asset dependency and no image file is redistributed with it. Licence: version arXiv:2509.24064v1 of this article is distributed under the Creative Commons Attribution 4.0 International licence, as recorded in the arXiv licence metadata for this exact version and shown on the abstract page https://arxiv.org/abs/2509.24064v1. A full rights notice is delivered at licenses/arxiv-2509.24064-CC-BY-4.0.txt. Characters: every retained excerpt is pure ASCII, so the delivered engine, XeLaTeX with its default Latin Modern text font, reports no missing character.
\begin{theorem}[Edge inducibility is ``as hard as'' vertex inducibility]\label{thm:eind-as-hard-as-ind}
    Let $G$ be a graph on $n$ vertices, and let $G'$ be obtained by attaching a pendant edge to each vertex of $G$. Then
\[ \frac{\ind(G)}{2^n n!}=\frac{ \eind(G')}{|\Aut(G')|} . \]
\end{theorem}

\begin{theorem}\label{thm:aind-eind-ind}
    Let $G$ be a graph and $M$ be a perfect matching in $G$ so that $M$ is a unique fractional perfect matching.
    Then 
    \[\aind(\ldg(G,M))\leq \frac{2^{\abs{M}}\abs{M}!}{\abs{\Aut(G)}}\eind(G)\leq \ind(\ldg(G,M)).\]
\end{theorem}

\begin{proof}
    We first prove the second half of the inequality.
    The proof is similar to that of \cref{thm:eind-as-hard-as-ind}.
    Suppose that $H$ is a graph with $m$ edges.
    We will construct a local digraph $H'$ with $m$ vertices such that every induced copy of $G$ in $H$ gives rise to many distinct induced copies of $\ldg(G,M)$ in $H'$. This is done as follows.
    For each $e\in E(H)$, let $v_e$ be a vertex in $H'$ and if $e=\{u,v\}$, fix any bijection $\sgn_e\colon\{u,v\}\to \{\pm\}$.
    For any edges $e,e_1,e_2\in E(H)$ where $e_1,e_2$ are disjoint and $e$ goes between $e_1$ and $e_2$, we create an edge $\Tilde{e}$ between $v_{e_1},v_{e_2}$. 
    We also set
    \[\sgn(v_{e_i},\Tilde{e}) = \sgn_{e_i}(e_i\cap e)\]
    for $i\in\{1,2\}$.
     From this construction and the definition of $\ldg(G,M)$, in any induced copy of $G$ in $H$, the vertices corresponding to the unique perfect matching of the induced copy induce a local digraph isomorphic to $\ldg(G,M)$ in $H'$.
    Therefore, we know that 
    \[\frac{\abs{M}!\cdot N_{\ind}(\ldg(G,M),H')}{m^{\abs{M}}}\geq \frac{\abs{M}!\cdot N_{\ind}(G,H)}{m^{\abs{M}}}= \frac{2^{\abs{M}}\abs{M}!}{\abs{\Aut(G)}}\cdot\frac{ \abs{\Aut(G)}N_{\ind}(G,H)}{(2m)^{\abs{M}}}.\]
    By the definitions of edge inducibility for graphs and vertex inducibility for local digraphs, we get that
    \[\ind(\ldg(G,M))\geq  \frac{2^{\abs{M}}\abs{M}!}{\abs{\Aut(G)}}\eind(G),\]
    as desired.

    To show the first half of the inequality, we begin with an LDAG $H$ with $n$ vertices.
    We will construct a graph $H'$ with many induced copies of $G$ corresponding to indueced copies of $\ldg(G,M)$ in $H$, though the number of edges of $H'$ would not be $n$.
    Instead, we will construct a graph $H'$ with $(n+o_{n;k\to\infty}(1))k$ edges where $k\in\NN$ is a sufficiently large positive integer.
    To do so, let $v_1,\ldots, v_n$ be a topological sort of $H$, and also relabel the signs so that $\sgn(v_i,e)=+$ for any $e$ going between $v_i,v_j$ and $i<j$.
    Take the graphification $\gr(H)$ of $H$, and let $M_H$ be the perfect matching $\{v_1^+v_1^-,\ldots, v_n^+v_n^-\}$.
    Then any induced copy of $\ldg(G,M)$ in $H$ corresponds to an induced copy of $G$ in $\gr(H)\cup M_H$ where the unique matching $M$ of $G$ is mapped to a submatching of $M_H$.
    We can also, without losing any induced copies of $G$, assume that $\gr(H)\cup M_H$ is simple.
    The graph $H'$ will then be chosen to be an appropriate blowup of $\gr(H)\cup M_H$ constructed as follows.
    Fix some rational numbers $0<\alpha_1<\cdots <\alpha_n<1/2$, and let $k$ be a large positive integer such that $k^{\alpha_i}$ is also a positive integer for any $i\in[n]$.
    Let $H'$ be the blowup of $\gr(H)\cup M_H$ where $v_i^+$ is replaced by an independent set of size $k^{\alpha_i}$ and $v_i^-$ is replaced by an independent set of size $k^{1-\alpha_i}$.
    Then the number of edges in $H'$ is at least $nk$ and can be upper bounded by
    \[nk+\sum_{1\leq i<j\leq n}k^{\alpha_i}\left(k^{\alpha_j}+k^{1-\alpha_j}\right) = (n+o_{n;k\to\infty}(1))k.\]
    This shows that $e(H') = (n+o_{n;k\to\infty}(1))k$.
    Moreover, for any induced copy of $G$ in $\gr(H)\cup M_H$ whose unique perfect matching $M$ is a subset of $M_H$, it is blown-up to $k^{\abs{M}}$ induced copies in $H'$.
    Therefore,
    \[ N_{\ind}(\ldg(G,M),H)\leq  \frac{N_{\ind}(G,H')}{k^{\abs{M}}}
        =(1+o_{n,\abs{M};k\to\infty}(1))(2n)^{\abs{M}}\cdot \frac{N_{\ind}(G,H')}{(2e(H'))^{\abs{M}}}.\]
    By taking $k$ to infinity, we see that
    \[N_{\ind}(\ldg(G,M),H)\leq \frac{(2n)^{\abs{M}}}{\abs{\Aut(G)}}\eind(G),\]
    and so
     \[\frac{\abs{M}!\cdot N_{\ind}(\ldg(G,M),H)}{n^{\abs{M}}}\leq \frac{2^{\abs{M}}\abs{M}!}{\abs{\Aut(G)}}\eind(G).\]
    By the definition of acyclic vertex inducibility, we see that $\aind(\ldg(G,M))\leq \frac{2^{\abs{M}}\abs{M}!}{\abs{\Aut(G)}}\eind(G)$, as desired.
\end{proof}

\end{document}
